\documentclass{article}
\usepackage[margin=2cm]{geometry}
\usepackage{amsmath}
\usepackage{graphicx}
\usepackage{booktabs}
\usepackage[utf8]{inputenc}
\usepackage{ragged2e}
\setlength{\emergencystretch}{3em}
\begin{document}
\sloppy

\noindent\hfill\includegraphics[height=1.8em]{assets/figure_header_logo.png}\par\noindent\hrule\vspace{4pt}

We can write


\begin{equation}
\begin{array} { r l } & { \mathfrak { J } _ { 3 _ { D } } ^ { 3 _ { B } ~ \big ( 1, 0 ; ( \tilde { A } \tilde { B } / \tilde { D } \tilde { C } ) ^ { \frac { 1 } { 2 } } \big ) } \mathfrak { F } ( u ) = \frac { \prod _ { i = 1 } ^ { 4 } ( 1 - a _ { i } u ) } { \big ( 1 - u ^ { 2 } \big ) \big ( 1 - q u ^ { 2 } \big ) } ( \mathfrak { F } ( q u ) - \mathfrak { F } ( u ) ) } \\ & { \qquad + \frac { \prod}}\end{array}
\end{equation}


Thefirsttwotermsdefine therank oneKoornwinder operator andwehave
an additional constant term. The eigenfunctions are the Askey-Wilson polynomials.In general these polynomials have four independent parameters ai
but in our case they are restricted to aia4 = a2a3. The constant term is


\begin{equation}
\begin{array} { r l } & { E _ { 3 _ { D } } ^ { 3 n, \left ( \tilde { A } \widetilde { B } / \tilde { D } \widetilde { C } \right ) ^ { \frac { 1 } { 2 } } \right ) } = \left [ - \frac { q ^ { 4 } t ^ { 2 } } { ( \tilde { A } \widetilde { C } ) ^ { 2 } } - q ^ { 2 } t ^ { 2 } \left ( 1 - \frac { \widetilde { B } } { \widetilde { C } \widetilde { D } \widetilde { A } } - \frac { 1 } { \widetilde { C } ^ { 2 } } - \frac { 1 + t ^ { - 2}}}\end{array}\right.
\end{equation}


Let us take the limit of the three punctured sphere.We give details of the
computation in Appendix with the final result given here(taking lim of the
p 0
right-hand side as in (??))


\begin{equation}
\begin{array} { r l } & { T _ { 3 x, 3 C }, 3 _ { D } ( w, u, v ) = \frac { 1 } { \left ( q p \frac { 1 } { A B C D } t ^ { - 2 } ; q \right ) } \left ( q p ^ { 2 } \frac { 1 } { A B C D } ; q \right ) ^ { 2 } \left ( q p \frac { 1 } { B A C D } ; q \right ) \textbf { ( 1 ) } } \\ & { \times \frac { 1 } { \left ( \sqrt { q p } \frac { t } { A C } v ^ { \pm 1 }, \sqrt { q p } \frac {}}}\end{array}\right.
\end{equation}


The index factorizes. In particular on general grounds we expect the index
ofthethreepuncturedspheretobe


\begin{equation}
\begin{array} { r l } & { \frac { 1 } { \left ( \sqrt { q p } \frac { t } { A C } v ^ { \pm 1 }, \sqrt { q p } \frac { t } { D B } v ^ { \pm 1 } ; q \right ) } \frac { 1 } { \left ( \sqrt { q p } \frac { t } { A B } w ^ { \pm 1 }, \sqrt { q p } \frac { t } { D C } w ^ { \pm 1 } ; q \right ) } } \\ & { \quad \times ~ \frac { 1 } { \left ( \sqrt { q p } \frac { t } { B C } u ^ { \pm 1 }, \sqrt { q p } \frac}}\end{array}\right.
\end{equation}


where the sum is over all the eigenvalues of theKoornwinder operator and
yx(z) are Askey-Wilson polynomials.What we have shown above is that

\end{document}
